\section{Methods}\label{sec:methods}
\subsection{Data acquisition} UniProt flat file, QuickGO GOA rows, cached PubMed/PMC text,
Reactome, PANTHER/PAINT, InterPro. \todo{versions and dates.}
\subsection{Deep research} \todo{providers, prompts, how used.}
\subsection{Agentic review} \todo{agent, model versions, skills, tool access (OLS, PubMed, UniProt).}
\subsection{Validation} LinkML schema validation, term validation, verbatim quote validation,
best-practice and policy checks (\texttt{just validate}); compliance scoring.
\subsection{Residue-claim validation}\label{sec:methods-residue}
Residue claims are checked by \texttt{ai\_gene\_review.validation.gene\_residue\_claims} against UniProt
sequences at pinned sequence versions. Four checks run per claim: the anchor protein has the stated residue
at the stated position; the target has its stated residue, and its accession is the review's own protein;
\textsc{lost} requires the two residues to differ and \textsc{retained} requires them to match; and a cited
family site exists in that family review and declares the anchor position. Across the corpus the validator
reports \nResCheckPass{} passing checks, \nResCheckFail{} failures and \nResCheckUnresolved{} unresolved
(a claim with no target position). Matching was by multiple sequence alignment for \nResMethodMsa{} claims,
structural superposition for \nResMethodStructure{}, UniProt feature tables for \nResMethodUniprotFeature{}
and a publication's statement for \nResMethodLiterature{}. Only \nResPinned{} claims name the release of a
shared alignment resource; the rest come from alignments computed for the review, for which the two
sequence lookups, not the alignment, are the evidence. Residue numbers reported by papers on a
non-displayed isoform are translated to the displayed sequence and checked the same way
(\texttt{finding\_residue\_mappings}).
\todo{provide alignment files or hashes for the MSA-based claims in the data release.}

\subsection{Protein binding}\label{sec:methods-pb}
\go{GO:0005515}{protein binding} is the single most frequent term in the corpus (\nPB{} rows, \pctPB\%,
almost all IPI). It records that an interaction was observed, not what the protein does, and is
expected to be retired from GO \todo{cite}. Reviews remove it by policy, so we exclude it from all
summary statistics. The exception is informative: in \pbModify\% of these rows the reviewer could
\emph{deepen} the annotation to a molecular function the interaction evidence supports (adaptor,
scaffold, enzyme regulator, receptor ligand); see Table~\ref{tab:pb}.

\subsection{Corpus statistics} All numbers are produced by \texttt{paper/scripts/}
(\texttt{extract\_annotations.py}, \texttt{extract\_goa\_meta.py}, \texttt{compute\_stats.py}) from
the repository at commit \todo{hash}. Evidence classes: Experimental = EXP, IDA, IPI, IMP, IGI, IEP;
High-throughput = HTP, HDA, HMP, HGI, HEP; Phylogenetic = IBA, IBD, IRD, IKR; Sequence/orthology =
ISS, ISO, ISA, ISM, IGC, RCA; Author/curator = TAS, NAS, IC. ``Rejected'' = \act{REMOVE} or
\act{MARK\_AS\_OVER\_ANNOTATED}; ``changed'' adds \act{MODIFY} and \act{TOO\_GENERAL}.
\subsection{Data and code availability} \todo{repo, Zenodo DOI, website.}
